% Track r3 -- the field-strength sector X_EM = (B;E) on the Clifford octants.
% Verification: investigations/two_readings/verify/r3_octants.py (claims r3-1 ... r3-26).
\section{\texorpdfstring{The field-strength sector $X_{\rm EM}=(\bB;\bE)$ on the Clifford octants}{The field-strength sector X\_EM = (B;E) on the Clifford octants}}
\label{r3:sec:main}

The question asked in the introduction about the octants concerns the field reading $X_{\rm EM}=(\bB;\bE)$, where each of the six components
$E_x,E_y,E_z$, $B_x,B_y,B_z$ is to sit on an octant. In the sense of the working note \cite{FMB} (its Definition~2.2), $X_{\rm EM}$ is not a sum
of two vectors and not a displacement: the pair $(\bB;\bE)$ labels the six components of the field two-form $F$, and the question is how
these six components, as six independent components of one field strength, sit on the eight octants of the spinor module. Everything below
is a classical field representation: placing $(\bB;\bE)$ on six lines of $S$ quantises nothing. A construction of this kind was proposed in an earlier, superseded version of the programme: the eight octants are the joint
eigenspaces of the pole operators $\omega_i$; the six mixed octants are field octants, those with one plus sign (temporal weight one
there, Fock weight $N=2$ here) carrying the electric components (as boosts) and those with two plus signs ($N=1$) the magnetic
components (as rotations); the two poles carry ``mass'' and
``statistics''; and Maxwell's equations are ``pooled'' across the octants. This section establishes in the conventions of \cite{ES17}
what of this is true. In brief: the octants are the eight lines of a Fock basis of the spinor module $S=\R^8$, which is the exterior
algebra of a null three-plane; there is a field-octant map $\Phi$ placing the electric member $\bE$ of the pair on the Fock weight $N=1$ and the magnetic member $\bB$ on $N=2$, unique up to
a scale and a quarter turn of duality once rotations and Maxwell's equations are required, the quarter turn giving the earlier
labelling; it carries the rotations, a parity and the
premetric $\mathfrak{gl}(3)$, but no even element of $\Cl(3,3)$ carries the boosts or the duality rotation; the Spin$(3,3)$-invariant mass
form pairs each octant with its antipode and $\bE$ with $\bB$; and Maxwell's four equations are one first-order equation on $S$, with the
two Gauss laws on the poles. In textbook variables the whole construction is Hestenes' Pauli-algebra electrodynamics, written with
the three temporal generators acting on the vacuum pole.

\paragraph{Conventions.} $V=\R^{3,3}$, $\eta=\operatorname{diag}(1,1,1,-1,-1,-1)$, the generators $e_1,\dots,e_6$ are those of eq.~(1) of
\cite{ES17}, $\omega:=e_1\cdots e_6=\Gamma$, $T:=e_1e_2e_3$ (the mass form, Corollary~18.4 of \cite{ES17}), and $\mathsf P=\sum_i\omega_i$,
$\omega_i:=e_ie_{i+3}$. The field point is $X=\sum_i(B_ie_i+E_ie_{i+3})$ (Theorem~12.3 of \cite{ES17}), with $\bB$ minus the textbook magnetic
field, $\Btext=-\bB$; $z=\bB+i\bE$, $q=z\cdot z$, $E_6=\tfrac12|z|^2$, $E_0=\tfrac12|q|$ (eq.~(1) of \cite{RE}). The binding root is
$Ge_i=e_{i+3}$, $Ge_{i+3}=-e_i$; the grading root of \cite{ES17} is $\cK=\tfrac12\ad_{\mathsf P}$, $\cK e_i=-e_{i+3}$, $\cK e_{i+3}=-e_i$
(Proposition~3.3 and Lemma~3.4 of \cite{ES17}; the para-complex structure of the earlier version is $-\cK$). On $V$ we use coordinates
$(\bB;\bE)\in\R^6$. Units are Heaviside--Lorentz with $c=1$ unless SI is stated. Dimensions: in SI $[\bE]=\mathrm{V\,m^{-1}}$,
$[\bB]=\mathrm T$, so that $\bE$ and $c\bB$ share a unit; $[t]=\mathrm s$, $[x^a]=\mathrm m$, $[\partial_a]=\mathrm m^{-1}$; $[\rho]=\mathrm{C\,m^{-3}}$,
$[\bj]=\mathrm{A\,m^{-2}}$, $[\varepsilon_0]=\mathrm{F\,m^{-1}}$; the generators, the Fock operators, the octant basis vectors and $T$ are
dimensionless, and $\psi=\Phi(X)$ carries the dimension of a field strength. Identifiers r3-$n$ refer to the script
\texttt{verify/r3\_octants.py}, which checks every finite-dimensional and symbolic claim below exactly.

\subsection{The octants and the octant group}
\label{r3:sec:octants}

\begin{proposition}[Octants and Fock grading]
\label{r3:prop:octants}
\begin{enumerate}[label=(\roman*),leftmargin=2.2em]
\item The pole operators are $\omega_i|b\rangle=(-1)^{b_i}|b\rangle$ on the basis $|b_1b_2b_3\rangle$ of $S$; their joint eigenlines, the
\emph{octants}, are the eight basis lines, labelled by the sign triples $\sigma=(\sigma_1,\sigma_2,\sigma_3)$, $\sigma_i=(-1)^{b_i}$.
\item $a^\pm_i:=\tfrac12(e_i\pm e_{i+3})$ satisfy $\{a^+_i,a^+_j\}=\{a^-_i,a^-_j\}=0$, $\{a^+_i,a^-_j\}=\delta_{ij}$, and
$\omega_i=\id-2a^+_ia^-_i$. Hence $N:=\sum_ia^+_ia^-_i$ is diagonal and counts the minus signs of $\sigma$.
\item $|\mathrm{vac}\rangle:=|000\rangle$, of signs $(+,+,+)$, spans the common kernel of the $a^-_i$ (the Fock vacuum, $N=0$), and
$a^+_{i_1}\cdots a^+_{i_r}|\mathrm{vac}\rangle=|b_I\rangle$ for $i_1<\dots<i_r$, where $b_I$ has ones exactly at $I$. The filled pole is
$|\mathrm{full}\rangle:=|111\rangle$, of signs $(-,-,-)$ and $N=3$.
\item The octant $\sigma$ is the pure spinor line of the maximal isotropic plane $N_\sigma:=\operatorname{span}\{e_i-\sigma_ie_{i+3}\}$: the
vectors $v\in V$ with $v|b\rangle=0$ are exactly those of $N_\sigma$.
\item $\Gamma=-\omega_1\omega_2\omega_3$, so $(-1)^N=\omega_1\omega_2\omega_3=-\Gamma$, which is the chirality $\gamma_5$ of Section~5 of
\cite{ES17}; on the octant with $w=3-N$ plus signs $\Gamma=(-1)^w$.
\end{enumerate}
\end{proposition}

\begin{proof}
(i) is in Proposition~2.5 and Appendix~A of \cite{ES17}: $e_i|b\rangle=(-1)^{b_1+\dots+b_{i-1}}|b\oplus\delta_i\rangle$ and
$e_{i+3}|b\rangle=(-1)^{b_1+\dots+b_i}|b\oplus\delta_i\rangle$. (ii) The Clifford relations give $\{e_i\pm e_{i+3},e_j\pm e_{j+3}\}=2(\delta_{ij}-\delta_{ij})=0$
and $\{e_i+e_{i+3},e_j-e_{j+3}\}=4\delta_{ij}$; then $\omega_i=(a^+_i+a^-_i)(a^+_i-a^-_i)=a^-_ia^+_i-a^+_ia^-_i=\id-2a^+_ia^-_i$. (iii)/(iv) By the two
formulas of (i), $(e_i-s\,e_{i+3})|b\rangle=(-1)^{b_1+\dots+b_{i-1}}(1-s(-1)^{b_i})|b\oplus\delta_i\rangle$, which vanishes iff $s=\sigma_i$. So $N_\sigma$
annihilates the octant $\sigma$; an annihilator is isotropic, since $2\eta(u,v)\phi=(uv+vu)\phi=0$, hence of dimension at most three, and it equals
$N_\sigma$. For $\sigma=(+,+,+)$ this is the statement $a^-_i|000\rangle=0$. The sign of $a^+_{i_1}\cdots a^+_{i_r}|\mathrm{vac}\rangle$ is a product of
factors $(-1)^{b_1+\dots+b_{i-1}}$ taken while all bits before $i$ are still zero. (v) Reordering $e_1e_4e_2e_5e_3e_6$ into $e_1\cdots e_6$
costs three transpositions (Example~2.6 of \cite{ES17}).
\end{proof}

All five items are restated from \cite{ES17} or fixed here (r3-1, r3-2, r3-3). Item (v) confirms the earlier version's statement
$I=-\omega_1\omega_2\omega_3$ with eigenvalue $(-1)^w$; note only that its ``temporal weight'' $w$ counts plus signs, so $w=3-N$.

\begin{proposition}[The octant group]
\label{r3:prop:octgroup}
For each axis $a$, the automorphisms of $\Cl(3,3)$ extending the reflection $R_a:e_a\mapsto-e_a$, $e_b\mapsto e_b$ $(b\neq a)$ are the
conjugations by the multiples of $g_a:=e_a\omega$. Each $g_a$ is a signed permutation matrix; it conjugates $\omega_i$ to $-\omega_i$ for
$a\in\{i,i+3\}$ and fixes the other two, so it carries the octant $\sigma$ to the octant with $\sigma_i$ reversed. The reflections of a
temporal axis and of its spatial partner act alike on the octants, and $g_ig_{i+3}=\pm\omega_i$ fixes every octant. The lifts $g_1,g_2,g_3$
satisfy $g_i^2=-\id$, anticommute pairwise and generate a group of order $16$ whose centre $\{\pm\id,\pm g_1g_2g_3\}$ contains
$g_1g_2g_3=e_4e_5e_6$, the element mapping each octant to its antipode; the kernel of its action on the eight octants is $\{\pm\id\}$, and
it acts through $\Z_2^3$, simply transitively.
\end{proposition}

\begin{proof}
$\Cl(3,3)\cong M_8(\R)$ is central simple, so every automorphism is inner and its implementing element is unique up to a real factor
(Skolem--Noether). Since $\omega$ anticommutes with every generator, $g_ae_bg_a^{-1}=-e_ae_be_a^{-1}$, which is $-e_b$ for $b=a$ and $e_b$
otherwise. Hence $g_a\omega_ig_a^{-1}=R_a(e_i)R_a(e_{i+3})$. Further $g_i^2=-e_i^2\omega^2=-\id$, and $g_ig_j=-e_ie_j$ anticommutes with
$g_jg_i=-e_je_i$. Moving $g_i$ past the other two lifts costs two signs, so $g_1g_2g_3=-T\omega=e_4e_5e_6$ commutes with every $g_i$; it is
not $\pm\id$, and it reverses all three signs. The images of $|\mathrm{vac}\rangle$ under the eight sign classes are the eight octants
(r3-4, r3-5).
\end{proof}

Fixed. The earlier version defined the octant group as the reflections of the temporal $\R^3$ acting on its eight \emph{sign sectors}
$\{x\in V^+:\sigma_ix_i>0\}$, and identified them with the spinor octants. Both sets are $\Z_2^3$-torsors, so equivariant bijections
between them exist, eight of them, one for each choice of the octant assigned to the positive sector. One of them is singled out by an
operator of \cite{ES17}:

\begin{proposition}[Clock octant and vacuum]
\label{r3:prop:clock}
For a dimensionless $x\in V^+$ the conversion operator of Theorem~3.6 of \cite{ES17} with $C=\operatorname{diag}(x)$ is
$B_x=\sum_ix_ie_ie_{i+3}=\sum_ix_i\omega_i$. It has eigenvalue $\sum_i\sigma_ix_i$ on the octant $\sigma$; if no $x_i$ vanishes, its unique top
eigenline is the octant $\operatorname{sign}(x)$; and $g_iB_xg_i^{-1}=B_{R_ix}$. The Fisher--Rao octant of clocks of Section~3.3 of \cite{ES17},
the closed positive octant $x_i=2\sqrt{w_i}\ge0$ of the sphere of radius two, therefore corresponds to the Fock vacuum: for $x$ in its
interior the top eigenline of $B_x$ is $|\mathrm{vac}\rangle$, and for the symmetric clock $B_{\widehat W}=\mathsf P/\sqrt3$ has top eigenvalue
$\sqrt3$ on $|\mathrm{vac}\rangle$. The weights $w_i=x_i^2/4$ do not see the signs of the amplitudes: the eight octants of the clock sphere are
the eight lifts of the clock simplex, which is their quotient by the octant group.
\end{proposition}

\begin{proof}
$B_x$ is diagonal by \cref{r3:prop:octants}(i), and $\sum_i\sigma_ix_i<\sum_i|x_i|$ unless $\sigma=\operatorname{sign}(x)$. Equivariance follows
from \cref{r3:prop:octgroup} (r3-6, with twenty random rational $x$).
\end{proof}

The spectral statements are fixed. The identification of the clock octant with the vacuum pole is constructed: it uses the order in
$\omega_i=e_ie_{i+3}$ and the choice of the top rather than the bottom eigenline, and the opposite choice assigns the filled pole. The two
objects are different in kind, a two-dimensional region of the sphere of clock directions in $V^+$ and a line of the spinor module; read
physically, the correspondence is interpretive.

\subsection{The exterior-algebra structure}
\label{r3:sec:exterior}

\begin{proposition}[The spinor module as an exterior algebra]
\label{r3:prop:fock}
$L:=\operatorname{span}\{e_i+e_{i+3}\}$ and $L':=\operatorname{span}\{e_i-e_{i+3}\}$ are maximal isotropic, transverse, and dual to each other
under $\eta$. The map $\Lambda^r(L)\to S$, $\ell_1\wedge\dots\wedge\ell_r\mapsto2^{-r}\ell_1\cdots\ell_r|\mathrm{vac}\rangle$, is an isomorphism
$\Lambda^*(L)\cong S$ carrying the degree to $N$, with $\dim\Lambda^w=1,3,3,1$ for $w=N=0,1,2,3$: the weights of the octants. Writing
$X=X_L+X_{L'}$, a vector acts as exterior multiplication by $2X_L$ plus contraction with the form $\eta(X_{L'},\cdot)$ on $L$,
\begin{equation}
\label{r3:eq:Xfock}
X=\sum_i\bigl[(B_i+E_i)\,a^+_i+(B_i-E_i)\,a^-_i\bigr].
\end{equation}
Moreover $a\mapsto a|\mathrm{vac}\rangle$ is an isomorphism from the temporal Clifford algebra $\Cl(V^+)=\langle e_1,e_2,e_3\rangle\cong\Cl(3,0)$ onto
$S$, and $e_I|\mathrm{vac}\rangle=a^+_I|\mathrm{vac}\rangle$: the octant weight $N$ is the Clifford grade of the temporal blade $e_I$.
\end{proposition}

\begin{proof}
$\eta(e_i\pm e_{i+3},e_j\pm e_{j+3})=\delta_{ij}-\delta_{ij}=0$, while $\eta(e_i+e_{i+3},e_j-e_{j+3})=2\delta_{ij}$. Elements of $L$ anticommute, so the
map is well defined, and it is bijective by \cref{r3:prop:octants}(iii). Under it $a^+_i$ is $\ell_i\wedge$ with $\ell_i=e_i+e_{i+3}=2a^+_i$, and
$\ell'_i=e_i-e_{i+3}=2a^-_i$ is the contraction with $\eta(\ell'_i,\cdot)$, since $\eta(\ell'_i,\ell_j)=2\delta_{ij}$. \Cref{r3:eq:Xfock} is
$e_i=a^+_i+a^-_i$, $e_{i+3}=a^+_i-a^-_i$. Finally
$e_j=a^+_j+a^-_j$ and $a^-_j$ annihilates $a^+_I|\mathrm{vac}\rangle$ for $j\notin I$, which gives $e_I|\mathrm{vac}\rangle=a^+_I|\mathrm{vac}\rangle$ by induction (r3-7).
\end{proof}

Inherited: this is Chevalley's realisation of the spinors of a space $L\oplus L'$ with $L'\cong L^*$ on $\Lambda^*(L)$
\cite{Chevalley54,LawsonMichelsohn}, and the pure spinor lines of \cref{r3:prop:octants}(iv) are the lines of Cartan and Chevalley (see Section~2
of \cite{Gualtieri03}). The last statement is the classical identification of the spinors of $\Cl(3,3)\cong M_8(\R)$ restricted to
$\Cl(3,0)$ with the Pauli algebra itself.

\begin{proposition}[The commutant of the grading root]
\label{r3:prop:gl3}
The elements of $\so{3,3}$ commuting with $\cK$ form a subalgebra of dimension nine, the adjoint image of the bivectors
$E_{ij}:=a^+_ia^-_j-\tfrac12\delta_{ij}\id$ (so $E_{ij}=a^+_ia^-_j$ for $i\neq j$ and $E_{ii}=-\tfrac12\omega_i$). They commute with $N$; $\ad E_{ij}$ acts
on $L$ by the matrix unit $e_{ij}$ and on $L'$ by $-e_{ji}$, so the subalgebra is $\mathfrak{gl}(L)\cong\mathfrak{gl}(3,\R)$; and the spin lift of
$A\in\mathfrak{gl}(3)$ acts on $\Lambda^w(L)$ as $\Lambda^w(A)-\tfrac12\tr A$. As a $\mathfrak{gl}(3)$-module, $S\cong\Lambda^*(L)\otimes(\det L)^{-1/2}$.
\end{proposition}

\begin{proof}
$\cK$ has the eigenspaces $L$ (eigenvalue $-1$) and $L'$ ($+1$), so an element commuting with it preserves both; $\eta$-skewness then
forces $A|_{L'}=-(A|_L)^*$ under $L'\cong L^*$, and every $A|_L\in\mathfrak{gl}(L)$ occurs: dimension nine. Since $\cK=\tfrac12\ad_{\mathsf P}$ and
$\mathsf P=3\,\id-2N$, an element $\ad\beta$ commutes with $\cK$ iff $[N,\beta]=0$. A bivector is quadratic in the $a^\pm$; the number-conserving
ones are the $a^+_ia^-_j$ up to scalars, and $E_{ij}$ removes the scalar part. By the CAR $[a^+_ia^-_j,a^+_k]=\delta_{jk}a^+_i$ and
$[a^+_ia^-_j,a^-_k]=-\delta_{ik}a^-_j$, and $a^+_ia^-_j$ acts on $a^+_K|\mathrm{vac}\rangle$ as the derivation replacing $a^+_j$ by $a^+_i$ (r3-8).
\end{proof}

The statement is standard (Example~2.12 of \cite{Gualtieri03} gives $\Lambda^*V^*\otimes(\det V)^{1/2}$ with $V^*$ in the role of $L$); it is fixed
here in the realisation. It shows that $\mathfrak{gl}(3)$ does not act on $\Lambda^w$ as the bare exterior power: the density factor
$-\tfrac12\tr A$, the same on all weights, is part of the spin representation. The commutant of $\cK$ is not stated in \cite{ES17} or \cite{RE}.

\subsection{The field-octant map}
\label{r3:sec:phi}

The map should intertwine the rotations of the field. Which rotations act on both triples as the field's rotations is fixed first.

\begin{lemma}[The diagonal rotations]
\label{r3:lem:rot}
For cyclic $(i,j,k)$ the bivector $L_i=\tfrac12(e_{j+3}e_{k+3}-e_je_k)$ of \cite{RE} acts on $V$ by $\delta\bB=e_i\times\bB$, $\delta\bE=e_i\times\bE$
and commutes with $G$ and $\cK$. The paired bivector $e_je_k+e_{j+3}e_{k+3}$ rotates the two triples in opposite senses and commutes with
neither. The diagonal rotations $\so3_\Delta=\operatorname{span}\{L_i\}$ are the compact part of $\lG$ and the antisymmetric part of
$\mathfrak{gl}(3)$, and $\so3_\Delta=\lG\cap\mathfrak{gl}(3)$. In Fock form
\[
L_i=a^+_ka^-_j-a^+_ja^-_k,\qquad K_i=\tfrac12(e_ke_{j+3}-e_je_{k+3})=a^-_ja^-_k-a^+_ja^+_k .
\]
The spin lift of $L_i$ (left multiplication) satisfies $L_i^3=-L_i$, so $\exp(2\pi L_i)=\id$ on $S$; it annihilates exactly
$\operatorname{span}\{|\mathrm{vac}\rangle,|\mathrm{full}\rangle\}$ and acts as $L_if_m=\epsilon_{imn}f_n$, $L_ig_m=\epsilon_{imn}g_n$ on
\[
f_k:=a^+_k|\mathrm{vac}\rangle,\qquad g_m:=\tfrac12\epsilon_{mjk}\,a^+_ja^+_k|\mathrm{vac}\rangle .
\]
\end{lemma}

\begin{proof}
$[e_je_k,e_k]=2e_j$ and $[e_je_k,e_j]=-2e_k$, while $[e_{j+3}e_{k+3},e_{k+3}]=-2e_{j+3}$ because $e_{k+3}^2=-\id$: the relative minus sign makes the
triples turn alike. $G$ and $\cK$ map $e_i$ to $\pm e_{i+3}$, so they commute with $R\oplus R'$ iff $R=R'$. The Fock forms follow from
$e=a^++a^-$, $e'=a^+-a^-$ and the CAR, and the action on $f,g$ from $[L_i,a^+_m]=\epsilon_{imn}a^+_n$ and $L_i|\mathrm{vac}\rangle=0$. The intersection is
computed in r3-9, which also checks $K_i=GL_i$ on $V$ (Proposition~13.6 of \cite{ES17}) and the brackets of Theorem~12.3 of \cite{ES17}.
\end{proof}

Restated and fixed (r3-9). The paired combination $e_je_k+e_{j+3}e_{k+3}$, a natural first guess, is the wrong one; with the minus sign it is the diagonal
rotation of Theorem~12.1 of \cite{ES17}, whose spin lift has integer spin on $S$, as Proposition~13.12 of \cite{ES17} states.

\begin{theorem}[The field-octant map]
\label{r3:thm:phi}
\begin{enumerate}[label=(\roman*),leftmargin=2.2em]
\item The linear maps $V\to S$ intertwining $\so3_\Delta$ form a four-dimensional space. Those sending $\bE$ into weight one and $\bB$ into
weight two form the two-parameter family
\[
\Phi_{\alpha,\beta}(X)=\alpha\sum_kE_kf_k+\beta\sum_mB_mg_m ,
\]
and $\Phi:=\Phi_{1,1}$ is an isometry of $\R^6$ onto the six field octants.
\item Under $a^+_I|\mathrm{vac}\rangle\leftrightarrow dx^I$, $\Phi(X)=E_k\,dx^k+\tfrac12\epsilon_{mjk}B_m\,dx^j\wedge dx^k$: $\bE$ is a one-form and $\bB$ a two-form
on $\R^3$.
\item The spin lift of $\mathfrak{gl}(3)$ preserves the image of $\Phi$ and induces $\bE\mapsto(A-\tfrac12\tr A)\bE$ and $\bB\mapsto(\tfrac12\tr A-A^{\!\top})\bB$: $\bE$
transforms in $L\otimes(\det L)^{-1/2}$ and $\bB$ in $L^*\otimes(\det L)^{1/2}$, contragredient up to the determinant.
\item The same $\mathfrak{gl}(3)\subset\so{3,3}$ acts on the field point through the light-cone components, $\bB+\bE\mapsto A(\bB+\bE)$,
$\bB-\bE\mapsto-A^{\!\top}(\bB-\bE)$. For symmetric $A$ this is $\bB\mapsto A\bE$, $\bE\mapsto A\bB$, a set of boosts, and $\Phi$ intertwines the
vector action of $\mathfrak{gl}(3)$ with its spin lift only on $\so3_\Delta$.
\end{enumerate}
\end{theorem}

\begin{proof}
(i) As $\so3_\Delta$-modules $V=2\cdot\R^3$ and, by \cref{r3:lem:rot}, $S=2\cdot\R^3\oplus2\cdot\R$ (the $f$'s, the $g$'s and the poles). The
vector representation is absolutely irreducible, so by Schur's lemma \cite{GoodmanWallach} $\Hom(\R^3,\R^3)=\R$ and $\Hom(\R^3,\R)=0$; the
weight condition leaves one scale on each copy. (ii) $a^+_k$ acts as $dx^k\wedge$, and $g_1=a^+_2a^+_3|\mathrm{vac}\rangle\leftrightarrow dx^2\wedge dx^3$
cyclically. (iii) follows from \cref{r3:prop:gl3}: on $\Lambda^1$ the action is $A-\tfrac12\tr A$, and on $\Lambda^2(L)\cong L^*\otimes\det L$, in the
basis $g_m$, it is $\tr A-A^{\!\top}$, from which $\tfrac12\tr A$ is subtracted. (iv) is \cref{r3:prop:gl3} on $V$ with \cref{r3:eq:Xfock}
(r3-10, r3-11).
\end{proof}

Items (i), (iii), (iv) are fixed; $\Phi$ itself is constructed. Item (ii) is the premetric statement that the electric field strength is a
one-form and the magnetic field strength a two-form on space, $F=\bE\wedge d\sigma+\bB$ in the notation of Hehl and Obukhov \cite{HehlObukhov03}
(inherited). Item (iv) makes precise the statement that under $\mathfrak{gl}(3,\R)$ $\bE$ is a vector and $\bB$ a two-form: it is
true of the spin lift on the octants and false of the vector action on the field point, where the symmetric part of $\mathfrak{gl}(3)$
mixes $\bE$ with $\bB$. The two scales of (i) are tied together by Maxwell's equations in \cref{r3:prop:maxunique}.

\begin{proposition}[Parity on the octants]
\label{r3:prop:parity}
\begin{enumerate}[label=(\roman*),leftmargin=2.2em]
\item $\Ad_T|_V=\eta$ is the frame parity, fixing $\bB$ and reversing $\bE$ (Proposition~2.11 of \cite{RE}), and $T$ is, up to a factor, the only
element of $\Cl(3,3)$ implementing it. $T$ carries each octant to its antipode,
\[
T|\mathrm{vac}\rangle=|\mathrm{full}\rangle,\quad Tf_k=g_k,\quad Tg_k=-f_k,\quad T|\mathrm{full}\rangle=-|\mathrm{vac}\rangle,
\]
and $T\Phi=-\Phi\circ G$: on the field octants the spinor lift of parity acts as minus the Hodge star, exchanging electric and magnetic octants.
\item $(-1)^N\Phi=\Phi\circ\eta$: the weight parity, which is the chirality $\gamma_5=-\Gamma$, acts as $-1$ on the image of $\bE$ and $+1$ on the image
of $\bB$; but $\Ad_{(-1)^N}=-\id$ on $V$.
\item The Dirac parity $P=\gamma^0=e_1$ of Theorem~5.3 of \cite{ES17} acts on the frame field by
$(\bB;\bE)\mapsto(B_1,-B_2,-B_3;-\bE)$, and $e_1f_1=|\mathrm{vac}\rangle$: it leaves the field octants.
\end{enumerate}
\end{proposition}

\begin{proof}
(i) $T$ commutes with $e_1,e_2,e_3$ and anticommutes with $e_4,e_5,e_6$; uniqueness is Skolem--Noether. $T=\prod_i(a^+_i+a^-_i)$ flips all three
signs, and the displayed values are r3-12. With them $T\Phi(X)=\sum_k(E_kg_k-B_kf_k)$, while $GX$ has $\bB'=-\bE$, $\bE'=\bB$. (ii) $(-1)^N$ is $-1$ on
$f$ and $+1$ on $g$; $\Gamma$ anticommutes with every vector. (iii) $e_1$ commutes with itself and anticommutes with the other generators,
and it flips the first sign (r3-12).
\end{proof}

Fixed. The polar and axial characters of $\bE$ and $\bB$ are present on the octants, but they are carried by the weight parity
$(-1)^N=-\Gamma$, whose Clifford action on $V$ is the total reflection $X\mapsto-X$, and not by the spinor lift $T$ of the parity, which acts on
the field octants as the duality rotation by $-\pi/2$. This is possible because $\Phi$ intertwines the rotations only, not the Clifford
action of $\mathrm{Pin}(3,3)$. The statement that parity acts as $-1$ on the electric and $+1$ on the magnetic octants is therefore false if
parity means the spinor lift of the spatial inversion, and true if it means the weight parity.

\begin{remark}[The earlier field-octant bijection]
\label{r3:rem:bijection}
The earlier version called the octants with one plus sign electric and those with two plus signs magnetic, the octant with its single
plus sign in pair $k$ carrying $E_k$ and the octant with its single minus sign in pair $m$ carrying $B_m$. Its electric octants are therefore
the lines of the $g_k$, of weight $N=2$, and its magnetic ones those of the $f_m$, of weight $N=1$: up to signs, which it does not fix, its
labelling is that of $\Phi\circ G$ (\cref{r3:prop:maxunique}), not of $\Phi$. In both labellings $E_k$ and $B_k$ sit on antipodal octants,
which the mass form pairs (\cref{r3:prop:energies}). The earlier bijection mapped octants to a basis of interior bivectors, a bijection of labels
and not a linear map; it is admissible as a labelling and equivalent to $\Phi$ by a duality rotation of a quarter turn, and nothing in the
octant structure prefers it. Its derivation of $E=cB$ from a count of ``temporal kernels'' is not reproduced: in the present construction the
ratio of the two scales of $\Phi_{\alpha,\beta}$ is $c$ in SI units, the conversion constant of Proposition~2.7 of \cite{RE}, fixed by Faraday's
law (\cref{r3:prop:maxunique} and \cref{r3:thm:maxwell}(iii)).
\end{remark}

\subsection{Clifford multiplication does not give the pair reading}
\label{r3:sec:poles}

The field point is also a Clifford vector, and one may ask what it does to the octants by left multiplication. The answer shows why the
pair reading $X_{\rm EM}=(\bB;\bE)$ needs the map $\Phi$ and is not obtained by multiplying.

\begin{proposition}[Light-cone amplitudes]
\label{r3:prop:poles}
For every $X=\sum_i(B_ie_i+E_ie_{i+3})$
\begin{equation}
\label{r3:eq:poles}
X|\mathrm{vac}\rangle=\sum_i(B_i+E_i)\,f_i,\qquad X|\mathrm{full}\rangle=\sum_i(B_i-E_i)\,g_i=\sum_i(-1)^{i-1}(B_i-E_i)\,|111\oplus\delta_i\rangle ,
\end{equation}
and $X^2=\eta(X,X)\id=\sum_i(B_i+E_i)(B_i-E_i)\,\id$. The amplitudes $B_i\pm E_i$ are the coordinates of $X$ on $L$ and $L'$, the eigen-coordinates of
$\cK$, with $\bB+\bE=\sqrt2\,\Re(e^{-i\pi/4}z)$ and $\bB-\bE=-\sqrt2\,\Im(e^{-i\pi/4}z)$. Further
\begin{equation}
\label{r3:eq:eighth}
X\bigl(|\mathrm{vac}\rangle+|\mathrm{full}\rangle\bigr)=\sqrt2\,\Phi\bigl(e^{\pi G/4}X\bigr).
\end{equation}
No spinor $\chi$ satisfies $X\chi=\Phi(X)$ for all $X$; the Clifford multiplications $X\mapsto X\chi$ that intertwine the rotations are those with
$\chi\in\operatorname{span}\{|\mathrm{vac}\rangle,|\mathrm{full}\rangle\}$, and none of them places $\bE$ and $\bB$ on separate weights.
\end{proposition}

\begin{proof}
\Cref{r3:eq:poles} follows from \cref{r3:eq:Xfock} with $a^-_i|\mathrm{vac}\rangle=0$, $a^+_i|\mathrm{full}\rangle=0$ and $a^-_i|\mathrm{full}\rangle=g_i$; $X^2$ follows from the
CAR. The duality rotation $e^{\pi G/4}$ gives $\bE'=(\bE+\bB)/\sqrt2$, $\bB'=(\bB-\bE)/\sqrt2$, which is \cref{r3:eq:eighth}. If $X\chi=\Phi(X)$, then
$e_1\chi=g_1$ and $e_4\chi=f_1$, i.e.\ $\chi=e_1|011\rangle=|111\rangle$ and $\chi=-e_4|100\rangle=|000\rangle$, a contradiction. $L_i(X\chi)=[L_i,X]\chi+XL_i\chi$, so
equivariance requires $L_i\chi=0$, and then $X(a|\mathrm{vac}\rangle+b|\mathrm{full}\rangle)$ puts $a(\bB+\bE)$ on weight one (r3-13).
\end{proof}

Fixed; the reading is interpretive. Read as a Clifford vector acting on the vacuum pole, the field point does not keep the two members of
the pair apart: it places on the $i$-th electric octant the sum $E_i+B_i$, which in textbook variables is $E_i-B_{{\rm t}i}$, and on the
filled pole it places $B_i-E_i=-(E_i+B_{{\rm t}i})$ on the $i$-th magnetic octant. These sums are the coordinates of $X$ on the null planes
$L$ and $L'$, the light-cone components of the para-complex structure; they mix the electric and the magnetic member and are not the
field strength. A literal reading ``$X=\bE+\bB$'' would therefore be wrong: the pair $(\bB;\bE)$ is placed on the octants by $\Phi$, which is
not Clifford multiplication by any spinor, and the Clifford action differs from $\Phi$ by a duality rotation of an eighth of a turn,
\cref{r3:eq:eighth}. (An eighth of a turn also separates the Klein correspondence from the frame reading in Theorem~13.13 of \cite{ES17}; the
maps are different and we draw no conclusion from the coincidence.)

\subsection{Boosts and duality cannot act by even elements}
\label{r3:sec:boosts}

\begin{theorem}[The even obstruction]
\label{r3:thm:noeven}
Every even element of $\Cl(3,3)$ commutes with $(-1)^N$ and maps an octant of weight $w$ into the weights $w$ and $w\pm2$; this includes $\so{3,3}$,
$\mathrm{Spin}(3,3)$ and the spin lifts of $\lG$ and $\lint$. No even $Q$ satisfies $Q\Phi=\Phi K_k$ for a boost $K_k$ of $\lG$, or $Q\Phi=\Phi G$.
At the level of groups, no element of $\Cl(3,3)$ of definite parity intertwines $e^{sK_k}$ $(s\neq0)$ with $\Phi$, nor $e^{\vartheta G}$ unless
$\vartheta\in\tfrac\pi2\Z$.
\end{theorem}

\begin{proof}
$(-1)^N=-\omega$ commutes with even and anticommutes with odd elements, and weights of equal parity differ by $0$ or $2$. By Example~12.4 of
\cite{ES17}, $K_3e_1=e_5$; so $\Phi(e_1)=g_1$ has weight two and $\Phi(K_3e_1)=f_2$ weight one, and an even $Q$ cannot map the first to the
second. Likewise $\Phi(e_1)=g_1$ and $\Phi(Ge_1)=\Phi(e_4)=f_1$. For the groups, $e^{sK_3}e_1=\cosh s\,e_1+\sinh s\,e_5$, so $\Phi(e^{sK_3}e_1)$ has
components of both parities when $s\neq0$, while an element of definite parity maps $g_1$ to a vector of definite parity; the same holds for
$e^{\vartheta G}e_1=\cos\vartheta\,e_1+\sin\vartheta\,e_4$ unless $\cos\vartheta\sin\vartheta=0$ (r3-14, r3-15).
\end{proof}

Obstructed, under the hypothesis that the operator has definite parity; the hypothesis covers every element of the spin group and of its
Lie algebra. The intertwiners that do exist are odd:

\begin{proposition}[Odd intertwiners and the octant representation]
\label{r3:prop:odd}
\begin{enumerate}[label=(\roman*),leftmargin=2.2em]
\item $-T\Phi=\Phi G$ and $Q_k\Phi=\Phi K_k$ for
\[
Q_k:=-TL_k=-\tfrac12e_k(\id-\omega_j\omega_l)=-\tfrac12(e_k+e_{k+3}\omega),\qquad\{k,j,l\}=\{1,2,3\},
\]
an odd element with parts of grades one and five. The solutions of $Q\Phi=\Phi K_k$ form an affine space of dimension $16$, of dimension $8$
among odd elements, and contain no Clifford vector; no Clifford vector implements $G$ on the octants.
\item On $S=\Cl(V^+)|\mathrm{vac}\rangle$ (\cref{r3:prop:fock}), $Q_k(a|\mathrm{vac}\rangle)=-\tfrac12[e_k,a]|\mathrm{vac}\rangle$ and $-T$ is left multiplication by the
central element $-e_1e_2e_3$ of $\Cl(V^+)$.
\item $\rho_{\rm oct}(L_k)=L_k$, $\rho_{\rm oct}(K_k)=Q_k$ is a representation of $\lG$ on $S$. Its Casimir $\sum_k\rho_{\rm oct}(K_k)^2-\rho_{\rm oct}(L_k)^2$
is $4$ on the field octants, the field-strength value, and $0$ on the poles, and its intertwiners $V\to S$ are $\operatorname{span}\{\Phi,\Phi G\}$.
\item The spin lift of $\lG$ has Casimir $3$ on all of $S$ and admits no nonzero intertwiner $V\to S$. Its boosts carry $f_k$ to $-|\mathrm{full}\rangle$
and $|\mathrm{vac}\rangle$ to $-g_k$: under the genuine Lorentz action the electric octants are the spatial components of the four-vector $S^+$
whose time component is the filled pole, and the magnetic octants those of the four-vector $S^-$ whose time component is the vacuum.
\end{enumerate}
\end{proposition}

\begin{proof}
(i) $T$ commutes with the same-sign bivectors, hence with $L_k$, so $-TL_k\Phi=-T\Phi L_k=\Phi GL_k=\Phi K_k$ by \cref{r3:prop:parity}(i) and $K_k=GL_k$.
The closed forms use $T(e_je_l)=-e_k$ for cyclic $(k,j,l)$ and $e_{k+3}\omega=-e_k\omega_j\omega_l$. The values of a solution on the two poles are free, eight
components each, four of them of the parity opposite to the pole (r3-15). (ii) $e_{k+3}\omega$ commutes with $e_1,e_2,e_3$ and sends $|\mathrm{vac}\rangle$ to $-e_k|\mathrm{vac}\rangle$,
so it acts on $a|\mathrm{vac}\rangle$ as $-a\,e_k$. (iii) $T$ commutes with $\rho_{\rm oct}(L)$ and $T^2=-\id$, so the brackets of $\lG$ follow from
those of $\so3$; $\rho_{\rm oct}(K)^2=-L^2$ gives the Casimir, and an intertwiner is a $T$-linear map $\C^3\to\C^3$ commuting with $\so3$, hence a
complex multiple of $\Phi$ (r3-16). (iv) The Casimir is Proposition~13.12 of \cite{ES17}; $S=2(\tfrac12,\tfrac12)$ contains no
$(1,0)\oplus(0,1)$; the boost formulas are $K_k=a^-_ja^-_l-a^+_ja^+_l$ on $f_k$ and $|\mathrm{vac}\rangle$ (r3-16).
\end{proof}

Constructed: the operators $Q_k$ and the representation $\rho_{\rm oct}$; fixed: the solution counts of (i) and the Casimir values and the
intertwiner space of (iii); restated and fixed: (iv). Item (ii) identifies the construction: in textbook variables (\cref{r3:thm:maxwell}) the
field octants carry Hestenes' Pauli-algebra field $F=\bE+I\Btext$ \cite{Hestenes66}, written with $\sigma_k=e_k$ and $I=T$, and the boosts and the
duality rotation act as they do in the Pauli algebra, through the commutator with $\sigma_k$ and through multiplication by $I$. The
commutator is a left multiplication combined with a right multiplication, and on $S$ the right multiplication is the odd element
$e_{k+3}\omega$: this is the reason the boosts are odd. Which symmetries of the field the octant picture carries is now settled.

\begin{corollary}[What the octants carry]
\label{r3:cor:carries}
By even elements of $\Cl(3,3)$ the field octants carry the rotations $\so3_\Delta$ and the premetric $\mathfrak{gl}(3)$ of \cref{r3:thm:phi}(iii),
through their spin lifts, and the parity, through $(-1)^N=-\omega$. They do not carry the boosts or the duality
rotations, which act only through the odd operators $Q_k$ and $-T$ of the representation $\rho_{\rm oct}$, the Pauli algebra of the temporal triple.
The spin-lifted Lorentz algebra $\lG$ makes the field octants parts of two four-vectors, not of a field strength.
\end{corollary}

The working note \cite{FMB} (Section~7) states that there is no nonzero linear Lorentz-equivariant map from the field-strength module to
the four-vectors. Item (iv) confirms the same fact for the spinor module, whose modules under $\lG$ are four-vectors (for $E_n$ itself
the statement is \cref{r2:prop:nolinear}); the note's statement that
$G$ can represent the Hodge star on the field (its eq.~(3) in Section~2 and its verdict table) is confirmed on the octants with a correction: the
operator representing it there is the odd element $-T$, the note's $I_T=e_1e_2e_3$ of its Section~5, not an element of $\mathrm{Spin}(3,3)$.

\subsection{Energies on the octants}
\label{r3:sec:energies}

\begin{proposition}[Energies and the Chevalley--Mukai pairing]
\label{r3:prop:energies}
Let $\psi=\Phi(X)=\psi_{\bE}+\psi_{\bB}$ with $\psi_{\bE}=\sum_kE_kf_k$, $\psi_{\bB}=\sum_mB_mg_m$.
\begin{enumerate}[label=(\roman*),leftmargin=2.2em]
\item $|\psi|^2=|\bE|^2+|\bB|^2=2E_6$ and $\psi^{\!\top}(-1)^N\psi=|\bB|^2-|\bE|^2=\Re q$.
\item $\psi_{\bE}^{\!\top}T\psi_{\bB}=-\bE\cdot\bB=-\tfrac12\Im q$, while $\psi^{\!\top}T\psi=0$ and $\psi^{\!\top}T\Gamma\psi=2\,\bE\cdot\bB=\Im q$. Hence
\[
E_0=\tfrac12|q|=\tfrac12\Bigl[\bigl(\psi^{\!\top}(-1)^N\psi\bigr)^2+4\bigl(\psi_{\bE}^{\!\top}T\psi_{\bB}\bigr)^2\Bigr]^{1/2}.
\]
\item The $\mathrm{Spin}(3,3)$-invariant bilinear forms on $S$ are $\operatorname{span}\{T,T\Gamma\}$, $T$ antisymmetric and $T\Gamma$ symmetric. Under
$\Lambda^*(L)\cong S$ the Chevalley--Mukai pairing $(\varphi,\chi)=[\tau(\varphi)\wedge\chi]_{\rm top}$, with $\tau$ the reversal, equals
$-\varphi^{\!\top}T\chi$. Two octants pair to a nonzero value iff their null planes $N_\sigma,N_\tau$ are transverse, iff they are antipodal.
\item $(-1)^N$ is invariant under the six compact bivectors and under $\rho_{\rm oct}$, and under none of the nine boosts. The Euclidean norm is
invariant under the compact bivectors and the $\rho_{\rm oct}(L_k)$, but not under the boosts $\rho_{\rm oct}(K_k)=Q_k$, which are symmetric:
$|\psi|^2=2E_6$ changes under boosts, $|\Phi(e^{sK_3}e_1)|^2=\cosh2s$.
\item In SI units, with $\psi=\Phi(c\bB;\bE)$ $[\mathrm{V\,m^{-1}}]$, $u=\tfrac{\varepsilon_0}2|\psi|^2$ is the energy density and
$\psi^{\!\top}(-1)^N\psi=\Re q_c$, $q_c=z_c\cdot z_c$, $z_c=c\bB+i\bE$ (Proposition~2.8 of \cite{RE}).
\end{enumerate}
\end{proposition}

\begin{proof}
The $f_k,g_m$ are orthonormal, $(-1)^N=\mp1$ on them, $Tg_m=-f_m$ and $\Gamma=-(-1)^N$, which gives (i), (ii) and the formula for $E_0$ by
$|q|^2=(\Re q)^2+(\Im q)^2$. (iii) A form $M$ is invariant iff $\beta^{\!\top}M+M\beta=0$ for the fifteen bivectors, a linear system with a
two-dimensional solution space (r3-18, the count of Theorem~13.13 of \cite{ES17}); $T$ commutes with the antisymmetric (same-sign) and
anticommutes with the symmetric (mixed) bivectors (Proposition~2.11 of \cite{RE}), and $\Gamma$, which anticommutes with every vector, commutes
with all of them. The Mukai identity is checked on all $64$ pairs of basis forms, and $\dim(N_\sigma\cap N_\tau)$ is the number of agreeing signs.
(iv) $e_k$ and $e_{k+3}\omega$ are symmetric matrices, so $Q_k^{\!\top}+Q_k=2Q_k\neq0$; (iv), (v): r3-17, r3-18.
\end{proof}

Items (i), (ii), (iv) are fixed, (iii) is inherited in substance and verified here, (v) restates \cite{RE}. By (i), the Maxwell
Lagrangian $L_M=-\tfrac12\Re q$ of \cref{r1:prop:maxwell} is $-\tfrac12\psi^{\!\top}(-1)^N\psi=\tfrac12(|\psi_{\bE}|^2-|\psi_{\bB}|^2)$ on the octants, half the
squared electric minus the squared magnetic octant amplitudes. The Chevalley form on the spinors
of $L\oplus L'$, antisymmetric for $\dim L=3$ since $(-1)^{3\cdot2/2}=-1$, and its non-vanishing on pure spinors exactly for transverse null
planes are Chevalley's \cite{Chevalley54}; Gualtieri identifies the form with the Mukai pairing (Section~2.4 of \cite{Gualtieri03}, whose
Proposition~2.21 restates the transversality criterion). So the Spin$(3,3)$-invariant mass form $T$ is, on the octants, the Chevalley--Mukai pairing of each octant with its
antipode: of the vacuum with the filled pole and of $f_k$ with $g_k$, that is, of $E_k$ with $B_k$, with the sign fixed by (ii):
$\psi_{\bE}^{\!\top}T\psi_{\bB}=-\bE\cdot\bB=-\tfrac12\Im q$. The Lorentz invariant $\Re q$ appears on
the octants through the form $(-1)^N$, which is not $\mathrm{Spin}(3,3)$-invariant, while the invariant symmetric form $T\Gamma$ gives $\Im q$;
both facts are consistent with \cref{r3:thm:noeven} and with the representation $\rho_{\rm oct}$, which preserves $(-1)^N$.

\subsection{Maxwell's equations on the octants}
\label{r3:sec:maxwell}

\begin{proposition}[Exterior derivative and codifferential]
\label{r3:prop:dk}
On $S$-valued fields put $d:=\sum_ia^+_i\partial_i$ and $\delta:=-\sum_ia^-_i\partial_i$. Then $d^2=\delta^2=0$,
\[
d-\delta=\sum_ie_i\partial_i=:D,\qquad d+\delta=\sum_ie_{i+3}\partial_i=:D',\qquad D^2=\Delta,\qquad D'^2=-\Delta,
\]
and under $f_k\leftrightarrow dx^k$, $g_m\leftrightarrow\star dx^m$, $|\mathrm{vac}\rangle\leftrightarrow1$, $|\mathrm{full}\rangle\leftrightarrow dx^1\wedge dx^2\wedge dx^3$, $d$ and $\delta$
are the exterior derivative and the codifferential of the Euclidean $\R^3$: $d(\bE\cdot f)=(\curl\bE)\cdot g$, $d(\bB\cdot g)=(\nabla\cdot\bB)|\mathrm{full}\rangle$,
$\delta(\bE\cdot f)=-(\nabla\cdot\bE)|\mathrm{vac}\rangle$, $\delta(\bB\cdot g)=(\curl\bB)\cdot f$.
\end{proposition}

\begin{proof}
$d^2=\sum_{ij}a^+_ia^+_j\partial_i\partial_j=0$ because $a^+_ia^+_j$ is antisymmetric, and likewise for $\delta$; $a^+_i\pm a^-_i$ are $e_i$ and $e_{i+3}$, and
$D^2=\sum_ie_i^2\partial_i^2$, $D'^2=\sum_ie_{i+3}^2\partial_i^2$. The four formulas use $a^-_ig_m=\epsilon_{mik}f_k$ and $a^+_mg_m=|\mathrm{full}\rangle$ (r3-19).
\end{proof}

Inherited, and verified here: $D=d-\delta$ is K\"ahler's operator on inhomogeneous forms \cite{Kahler62}, the forms read as spinors \cite{Graf78}
and the Dirac--K\"ahler equation of Becher and Joos \cite{BecherJoos82}, here on space.

\begin{theorem}[Maxwell's equations as one equation on the octants]
\label{r3:thm:maxwell}
For every smooth field $X(t,\mathbf x)$,
\begin{equation}
\label{r3:eq:maxwell}
(\partial_t-D)\,\Phi(X)=-(\nabla\cdot\bE)\,|\mathrm{vac}\rangle+(\partial_t\bE+\curl\bB)\cdot f+(\partial_t\bB-\curl\bE)\cdot g-(\nabla\cdot\bB)\,|\mathrm{full}\rangle .
\end{equation}
Consequently:
\begin{enumerate}[label=(\roman*),leftmargin=2.2em]
\item Maxwell's equations with sources in the frame variables, $\partial_tz=-i\curl z-i\bj$ and $\nabla\cdot z=i\rho$ (eq.~(16) of \cite{ES17}), are equivalent to
\[
(\partial_t-D)\,\Phi(X)=-\bigl(\rho\,|\mathrm{vac}\rangle+\bj\cdot f\bigr):
\]
the vacuum component is Gauss's law $\nabla\cdot\bE=\rho$, the filled component the absence of magnetic charge, the electric octants carry the
Amp\`ere--Maxwell law $\partial_t\bE=-\curl\bB-\bj$ and the magnetic octants Faraday's law $\partial_t\bB=\curl\bE$.
\item $(\partial_t+D)(\partial_t-D)=\partial_t^2-\Delta$ on all eight components, while $(\partial_t+D')(\partial_t-D')=\partial_t^2+\Delta$. In $D=\sum_ie_i\partial_i$ the
derivatives $\partial_i$ are taken along the three physical spatial coordinates $x^i$, the lengths of the classical reading, which sit on $V^-$;
the Clifford generators multiplying them are the temporal $e_1,e_2,e_3$, of square $+\id$, the generators that carry $\bB$ in the field reading,
and not the spatial ones.
\item In SI units $\bigl(\tfrac1c\partial_t-D\bigr)\Phi(c\bB;\bE)=-(c\rho\,|\mathrm{vac}\rangle+\bj\cdot f)/(\varepsilon_0c)$.
\item In textbook variables $\Phi_{\rm t}(X):=\bE\cdot f+\Btext\cdot g=-(-1)^N\Phi(X)$ satisfies $T\Phi_{\rm t}=\Phi_{\rm t}\circ G$ and
\[
(\partial_t+D)\,\Phi_{\rm t}(X)=\rho\,|\mathrm{vac}\rangle-\bj\cdot f,\qquad \Phi_{\rm t}(X)=\Bigl(\sum_kE_ke_k+\sum_kB_{{\rm t}k}\,Te_k\Bigr)|\mathrm{vac}\rangle ,
\]
which, under $S\cong\Cl(V^+)$ with $\sigma_k=e_k$ and $I=T$, is the equation $(\partial_t+\nabla)(\bE+I\Btext)=\rho-\bj$ of the Pauli algebra.
\item With pole components $\psi=\Phi(X)+\phi_0|\mathrm{vac}\rangle+\phi_3|\mathrm{full}\rangle$, $(\partial_t-D)\psi=0$ reads $\partial_t\phi_0=\nabla\cdot\bE$, $\partial_t\phi_3=\nabla\cdot\bB$,
$\partial_t\bE=-\curl\bB+\nabla\phi_0$, $\partial_t\bB=\curl\bE+\nabla\phi_3$; every component of a solution obeys the wave equation, so the Gauss constraints
hold for all times if $\phi_0,\phi_3$ and the two divergences vanish initially.
\end{enumerate}
\end{theorem}

\begin{proof}
\Cref{r3:eq:maxwell} is $\partial_t\Phi(X)-(d-\delta)\Phi(X)$ with the four formulas of \cref{r3:prop:dk}; it is checked on general symbolic fields
$E_k(t,\mathbf x)$, $B_k(t,\mathbf x)$ (r3-20), and the real and imaginary parts of the residuals of eq.~(16) of \cite{ES17} are exactly the
components of $(\partial_t-D)\Phi(X)+\rho|\mathrm{vac}\rangle+\bj\cdot f$. (ii) $(\partial_t+D)(\partial_t-D)=\partial_t^2-D^2$ (r3-21). (iii) is (i) with $\Phi(\bE,c\bB)$ and
$\partial_t\bE=-c^2\curl\bB-\bj/\varepsilon_0$, $\nabla\cdot\bE=\rho/\varepsilon_0$ (r3-24, on first-order jets). (iv) $\Phi_{\rm t}(X)=\bE\cdot f-\bB\cdot g$; $(-1)^N$
anticommutes with $D$; $Te_k=e_je_l$ cyclically gives the second form; under $a|\mathrm{vac}\rangle\leftrightarrow a$ left multiplication by $e_k$ is that by $\sigma_k$ (r3-24). (v) is
\cref{r3:eq:maxwell} together with $D(\phi_0|\mathrm{vac}\rangle)=\nabla\phi_0\cdot f$ and $D(\phi_3|\mathrm{full}\rangle)=\nabla\phi_3\cdot g$ (r3-21); the wave equation follows from (ii), and
$\phi_0$ with $\phi_0=\partial_t\phi_0=0$ at $t=0$ vanishes identically.
\end{proof}

Fixed in the realisation and inherited in substance. Item (ii) is a statement about the representation and not about spacetime: the
octant form of Maxwell's equations uses the Pauli triple $e_1,e_2,e_3$ because it is the triple of square $+\id$ that makes the operator hyperbolic,
and this choice is constructed. It does not show that spatial differentiation follows from three temporal directions, and the spatial
coordinates $x^i$ remain the lengths of $V^-$; the Dirac operator of the particle reading couples the same derivatives to the spatial
generators through $\gamma^0\gamma^a$ instead (\cref{r4:sec:main}). Item (iv) is the Pauli-algebra form, obtained by multiplying with $\gamma_0$, of the
spacetime-algebra equation $\nabla F=J$ of Hestenes \cite{Hestenes66}, which Theorem~13.3 of \cite{ES17} realises on the interior four-plane; here
it is realised on the temporal triple acting on the vacuum. The element $T$ that plays the pseudoscalar $I$ is the $I_T$ of Section~5 of \cite{FMB},
and the properties stated there, $I_T^2=-\id$ and $[I_T,e_i]=0$ for the temporal $e_i$, are confirmed (r3-24). The circularly polarised wave of Example~12.11 of \cite{ES17} solves $(\partial_t-D)\Phi(X)=0$, and the same wave with $\bB$ reversed does
not (r3-23). The textbook signs, $\partial_t\bE=\curl\Btext-\bj$ and $\partial_t\Btext=-\curl\bE$, become those of (i) in the frame
variable $\bB=-\Btext$.

\begin{proposition}[Maxwell fixes the scales and leaves a duality angle]
\label{r3:prop:maxunique}
Let $\Phi_M(X)=(a\bE+b\bB)\cdot f+(c\bE+d\bB)\cdot g$ be the general rotation-equivariant map into the field octants. Then $(\partial_t-D)\Phi_M(X)$ vanishes
for every source-free Maxwell field iff $c=-b$ and $d=a$, that is iff $\Phi_M=r\,\Phi\circ e^{\vartheta G}$ for some $r\ge0$ and $\vartheta$. The map puts
$\bE$ and $\bB$ on single weights iff $\vartheta\in\tfrac\pi2\Z$: $\Phi$ for $\vartheta=0$ and $\Phi\circ G$, with $\bE$ on weight two, for $\vartheta=\tfrac\pi2$.
\end{proposition}

\begin{proof}
Substituting $\partial_t\bE=-\curl\bB$, $\partial_t\bB=\curl\bE$ and the vanishing divergences, the residual has $f$-component $(d-a)\curl\bB+(b+c)\curl\bE$,
$g$-component $(d-a)\curl\bE-(b+c)\curl\bB$ and no pole component; it vanishes on all first-order Maxwell jets iff $d=a$, $c=-b$ (r3-22). Conversely
$\Phi\circ e^{\vartheta G}$ maps solutions to solutions, since $e^{\vartheta G}$, the duality rotation $z\mapsto e^{i\vartheta}z$, preserves the source-free system.
\end{proof}

Fixed. The two scales of \cref{r3:thm:phi}(i) are tied by Faraday's and Amp\`ere's laws, $\beta=\alpha$ in Heaviside--Lorentz units and $\beta=c\,\alpha$
in SI, and the energy normalisation of \cref{r3:prop:energies}(i) fixes $\alpha$ up to sign. What remains is the duality angle. With electric
sources, $\Phi$ places the charge density on the vacuum pole and the current on the electric octants, and $\Phi\circ G$ places them on the
filled pole and on the weight-two octants; the convention that all charges are electric (Section~13 of \cite{ES17}, \cite{Jackson}) fixes which
octants carry the current once the labelling is chosen, not the labelling.

\begin{remark}[Octant pooling]
\label{r3:rem:pooling}
Each $e_i$ maps the octant $\sigma$ to the octant with $\sigma_i$ reversed, the image of $\sigma$ under the octant generator $R_i$, and
$e_{i+3}=e_i\omega_i$ acts on the same pair of octants (r3-25). Hence the $\sigma$-component of $D\psi$ collects $\partial_i$ of the three neighbouring
octants only: the Gauss law on the vacuum collects $\sum_i\partial_iE_i$ from the three adjacent electric octants, and each evolution equation
collects a curl from two octants of the other weight. This is the exact content of the proposed ``pooling'' of Maxwell's equations across
octants, and it holds as $D=\sum_ie_i\partial_i$. The proposed mechanism, a per-pair exchange of $e_i$ with $e_{i+3}$ converging on the pole
$(-,-,-)$, is not confirmed: the exchange replaces $D$ by $D'$ and the wave operator by the four-dimensional Laplacian
(\cref{r3:thm:maxwell}(ii)), and nothing converges on one pole, the two poles carrying one Gauss law each.
\end{remark}

\subsection{\texorpdfstring{The answer: $X_{\rm EM}=(\bB;\bE)$ with each component on an octant}{The answer: X\_EM = (B;E) with each component on an octant}}
\label{r3:sec:answer}

The pair reading has one exact realisation. The field-octant map (\cref{r3:thm:phi,r3:prop:maxunique}),
$\Phi(X)=\sum_k(E_k+B_kT)e_k|\mathrm{vac}\rangle$, puts each $E_k$ on the octant with a single minus sign in pair $k$ and each $B_k$ on its antipode; it
is unique up to a scale once rotation equivariance, Maxwell's equations and the separation of $\bE$ from $\bB$ are required, apart from the
alternative $\Phi\circ G$ that the earlier labelling uses; in textbook variables the same octants carry, through $\Phi_{\rm t}=-(-1)^N\Phi$
(\cref{r3:thm:maxwell}(iv)), Hestenes' $\bE+I\Btext$ on the temporal Pauli algebra, with $I=T$. Clifford multiplication by the field point does
not realise the pair: it places the light-cone sums $\bB\pm\bE$ on the octants (\cref{r3:prop:poles}), and differs from $\Phi$ by an eighth of a turn
of duality. The construction is a representation of the classical field strength; it is not a quantisation of the field.

\paragraph{The main obstruction.} On the octants the rotations act through the spin representation, but the boosts and the duality rotation
of the field act only through the odd operators $Q_k$ and $-T$ of $\rho_{\rm oct}$ (\cref{r3:thm:noeven,r3:prop:odd}), whereas a Dirac field transforms
under the spin lift of the interior Lorentz algebra. A covariant theory containing both the Dirac field and the field octants must say how
$\rho_{\rm oct}$, the spin lift acting on Dirac spinors and the map $\Phi$ coexist. This is the principal question the octant picture leaves, and the
element $T$, which is at once the mass form of the Dirac reading and the duality operator on the octants, is where the two meet
(\cref{r4:sec:main}).

\paragraph{What the poles carry.} With $\Phi$, the vacuum pole carries Gauss's law and the charge density, the filled pole the absence of
magnetic charge (\cref{r3:thm:maxwell}); they are the Fock vacuum and the filled state, the rotation scalars, the eigenlines of $\mathfrak{gl}(3)$ of
weights $\Lambda^0$ and $\Lambda^3$, and, under the spin-lifted $\lG$, the time components of the two four-vectors $S^\mp$ (\cref{r3:prop:odd}(iv)).
The mass form pairs them with each other as it pairs $E_k$ with $B_k$ (\cref{r3:prop:energies}). Of the earlier labels, ``statistics'' survives as
the constructed correspondence between the Fisher--Rao octant of clocks and the vacuum pole $(+,+,+)$ (\cref{r3:prop:clock}). ``Mass'' does not
survive as a property of the pole $(-,-,-)$: the mass form treats both poles alike, and the interior Lorentz algebra distinguishes no octant,
since the stabiliser in $\lint$ of the octant $\sigma$ is the three-dimensional span of $e_1e_4=\omega_1$ and the null rotations
$(e_1-\sigma_1e_4)e_5$, $(e_1-\sigma_1e_4)e_6$, depending on $\sigma_1$ only, while $g_2,g_3$ commute with $\lint$ and $g_1$ normalises it (r3-26). The
octant Fock structure is the canonical anticommutation algebra of one constituent; it is not a photon Fock space, and the quantisation of the
field stays open, as in \cite{FMB}.

\paragraph{Open.} Whether a dynamics, rather than a convention, selects $\Phi$ or $\Phi\circ G$; what the two scalar fields $\phi_0,\phi_3$ of the
extended system of \cref{r3:thm:maxwell}(v) mean, if anything; and how $\rho_{\rm oct}$, which is not part of the spin representation, enters a covariant
dynamics with the Dirac sector of \cite{ES17}; its covariant couplings linear in the field are determined in \cref{r4:thm:couplings}.

\Cref{r3:tab:claims} lists every claim of the section with its check and its class.

{\small\renewcommand{\arraystretch}{1.12}
\begin{longtable}{@{}>{\raggedright\arraybackslash}p{0.09\textwidth}>{\raggedright\arraybackslash}p{0.56\textwidth}>{\raggedright\arraybackslash}p{0.27\textwidth}@{}}
\caption{The claims of this section, with their checks and classes.}\label{r3:tab:claims}\\
\toprule
id & claim & class\\
\midrule
\endfirsthead
\toprule
id & claim & class\\
\midrule
\endhead
r3-1 & octants $=$ basis lines, sign triples, CAR, $N=\#$ minus signs, vacuum $(+,+,+)$ & restated (ES17), fixed\\
r3-2 & octant $\sigma$ is the pure spinor of $N_\sigma=\operatorname{span}\{e_i-\sigma_ie_{i+3}\}$ & fixed; inherited (Cartan--Chevalley)\\
r3-3 & $\Gamma=-\omega_1\omega_2\omega_3$, $(-1)^N=-\Gamma=\gamma_5$; the earlier $I=-\omega_1\omega_2\omega_3$ confirmed & restated, fixed\\
r3-4, 5 & reflections lift to $e_a\omega$, unique; octant group simply transitive; lifted group of order $16$, centre $\{\pm\id,\pm e_4e_5e_6\}$ & fixed\\
r3-6 & top eigenline of $B_{\operatorname{diag}x}$ is the octant $\operatorname{sign}x$; clock octant $\leftrightarrow$ vacuum & fixed; identification constructed\\
r3-7 & $S\cong\Lambda^*(L)\cong\Cl(V^+)|\mathrm{vac}\rangle$, weights $1,3,3,1$; $X$ acts by $2X_L\wedge$ and contraction with $\eta(X_{L'},\cdot)$ & inherited (Chevalley), verified\\
r3-8 & commutant of $\cK$ $=\mathfrak{gl}(3)$, acting as $\Lambda^w(A)-\tfrac12\tr A$ & inherited (Gualtieri), fixed here\\
r3-9 & diagonal rotations $L_i=\tfrac12(e_{j+3}e_{k+3}-e_je_k)$, Fock forms of $L_i,K_i$ & restated, fixed\\
r3-10 & equivariant maps: dimension $4$; $\Phi_{\alpha,\beta}$ & fixed; $\Phi$ constructed\\
r3-11 & premetric $\mathfrak{gl}(3)$ on the image; vector action mixes $\bE,\bB$ & fixed; premetric reading inherited\\
r3-12 & $T$ lifts parity and acts as $-G$ on the octants; $(-1)^N$ acts as parity & fixed\\
r3-13 & $X|\mathrm{vac}\rangle$, $X|\mathrm{full}\rangle$ light-cone amplitudes; \cref{r3:eq:eighth}; $\Phi$ is no Clifford map & fixed; reading interpretive\\
r3-14 & no even intertwiner of boosts or of $G$ & obstructed\\
r3-15, 16 & odd intertwiners $Q_k,-T$; $\rho_{\rm oct}$; spin lift: four-vectors, no intertwiner & $Q_k$, $\rho_{\rm oct}$ constructed; counts fixed; restated (iv)\\
r3-17, 18 & energies; invariant forms; $T=$ Chevalley--Mukai pairing of antipodes; $|\psi|^2$ not $\rho_{\rm oct}$-invariant & fixed; pairing inherited\\
r3-19 & $d,\delta$ on the octants, $d\mp\delta=D,D'$ & inherited (K\"ahler), verified\\
r3-20, 21, 23, 24 & Maxwell $\Leftrightarrow(\partial_t-D)\Phi(X)=-(\rho|\mathrm{vac}\rangle+\bj\cdot f)$; wave operator; SI; Hestenes form & fixed; inherited in substance\\
r3-22 & Maxwell-compatible equivariant maps $=r\Phi\circ e^{\vartheta G}$ & fixed\\
r3-25 & pooling $=$ nearest-neighbour structure of $D$ & fixed; mechanism refuted\\
r3-26 & $\lint$ distinguishes no octant & fixed; ``mass pole'' refuted\\
\bottomrule
\end{longtable}}
